% ---------------------------------------------------------------------------
% Where the two scales enter and how far apart the scales of the problem are.
% (a) the factorised picture with mu_F at the parton densities and mu_R at
% the vertices; (b) the scales of our event on a logarithmic ladder with the
% coupling from the -v log of standalone_qqbar_bbbar_xsec.py (section [2],
% step 3) and the value at mZ from the table.
% ---------------------------------------------------------------------------
\begin{tikzpicture}[x=1cm,y=1cm,
  every node/.style={font=\footnotesize},
  fermion/.style={thick,postaction={decorate},decoration={markings,mark=at position 0.55 with {\arrow[scale=1.1]{Stealth}}}},
  antifermion/.style={thick,postaction={decorate},decoration={markings,mark=at position 0.45 with {\arrowreversed[scale=1.1]{Stealth}}}},
  gluon/.style={thick,decorate,decoration={coil,aspect=0,segment length=4.5pt,amplitude=2pt}},
  pdfblob/.style={draw,circle,minimum size=9mm,thick,fill=blue!12},
  vertex/.style={fill=black,circle,inner sep=1.6pt},
  muF/.style={draw=blue!60!black,dashed,thick,rounded corners=3pt},
  muR/.style={draw=orange!80!black,dashed,thick,rounded corners=3pt},
  lab/.style={font=\scriptsize,align=center},
  head/.style={font=\small\bfseries},
  tick/.style={thick},
  ladder/.style={-{Stealth},very thick}]
% ---------------- (a) where mu_F and mu_R enter ----------------------------------------------------
\node[head,anchor=west] at (-0.3,3.0) {(a) two conventions, two scales};
\node[pdfblob] (FA) at (0.3,1.1) {$f_{\bar u}$};
\node[pdfblob] (FB) at (0.3,-1.1) {$f_{u}$};
\coordinate (V1) at (2.6,0); \coordinate (V2) at (4.6,0);
\draw[antifermion] (FA.east) -- (V1);
\draw[fermion] (FB.east) -- (V1);
\draw[gluon] (V1) -- (V2);
\draw[fermion] (V2) -- (6.4,1.3) node[right,lab] {$b$};
\draw[antifermion] (V2) -- (6.4,-1.3) node[right,lab] {$\bar b$};
\node[vertex] at (V1) {}; \node[vertex] at (V2) {};
\draw[muF] (-0.4,-1.75) rectangle (1.0,1.75);
\node[lab,text=blue!60!black,anchor=south] at (0.3,1.8) {$\muF$: what counts as\\``inside the proton''};
\draw[muR] (2.1,-0.55) rectangle (5.1,0.55);
\node[lab,text=orange!80!black,anchor=north] at (3.6,-0.6) {$\muR$: at which scale\\$\alphas$ is evaluated, twice};
\node[lab,text=black!65,text width=6.9cm,anchor=north west] at (-0.4,-2.35) {Radiation harder than $\muF$ belongs to $\hat\sigma$, softer radiation to the PDFs; the split is a convention and the full answer does not depend on it, but a leading-order answer does. \PYTHIA sets both to $\muR^2=\muF^2=\pthat^2+\mb^2$ (scale codes 1--3).};
% ---------------- (b) the ladder of scales ---------------------------------------------------------
\begin{scope}[xshift=9.6cm]
\node[head,anchor=west] at (-0.3,3.0) {(b) the scales of our event, $\log_{10}(\mu/\mathrm{GeV})$};
% axis from log10 = -0.7 (0.2 GeV) to 4.4 (25 TeV): 1 decade = 2.6 cm; x = (log10 v + 0.7) * 2.6
\draw[ladder] (0,0) -- (13.4,0);
\foreach \l/\t in {0/{$1$},1/{$10$},2/{$10^2$},3/{$10^3$},4/{$10^4$}} {\draw[tick] ({(\l+0.7)*2.6},-0.1) -- ({(\l+0.7)*2.6},0.1); \node[lab,below] at ({(\l+0.7)*2.6},-0.1) {\t};}
\draw[thick,gray!70] (0.46,0) -- (0.46,0.6);  \node[lab,anchor=south west,text=gray!80!black] at (0.3,0.6) {$\Lambda_{\mathrm{QCD}}\approx0.3\GeV$: perturbation theory fails};
\draw[thick,gray!70] (3.59,0) -- (3.59,1.4);  \node[lab,above,text=gray!80!black] at (3.59,1.4) {$\mb=4.8\GeV$\\threshold, $\rho=4\mb^2/\shat$};
\draw[thick] (6.92,0) -- (6.92,2.3); \node[lab,above] at (6.92,2.3) {$m_Z=91\GeV$: $\alphas=0.1180$, the input};
\draw[thick,orange!80!black] (7.16,0) -- (7.16,-0.8); \node[lab,below,text=orange!80!black] at (7.16,-0.8) {$\muR=\muF=Q=113\GeV$\\$\alphas(Q^2)=0.1143$};
\draw[thick,orange!60!black,dashed] (6.38,0) -- (6.38,-1.9); \node[lab,anchor=north east,text=orange!60!black] at (6.45,-1.9) {$Q/2=57\GeV$: $\alphas=0.1272$\\$\alphas^2$ up by $24\%$};
\draw[thick,orange!60!black,dashed] (7.94,0) -- (7.94,-1.9); \node[lab,anchor=north west,text=orange!60!black] at (7.87,-1.9) {$2Q=227\GeV$: $\alphas=0.1042$\\$\alphas^2$ down by $17\%$};
\draw[thick,red!70!black] (8.57,0) -- (8.57,1.0); \node[lab,anchor=south west,text=red!70!black] at (8.45,1.0) {$\sqrt{\shat}=396\GeV$: $\alphas=0.0977$\\(what scale code 4 would use)};
\draw[thick,gray!70] (12.57,0) -- (12.57,0.4); \node[lab,anchor=south east,text=gray!80!black] at (12.7,0.4) {$\sqrt{s}=13.6\TeV$};
\node[lab,text=black!65,text width=12.6cm,anchor=north west] at (-0.3,-3.1) {The hard process lives one decade above $\mb$ and four above $\Lambda_{\mathrm{QCD}}$. Moving $\muR$ by a factor two moves $\alphas^2$, and with it the whole leading-order rate, by about $20\%$: that is the uncertainty of \cref{sec:ancestral_home:scales}, and no amount of Monte Carlo statistics reduces it.};
\end{scope}
\end{tikzpicture}
